\documentclass[A4]{dnbqcm}
\usepackage{texmaths-tikz}
\usepackage{texmaths-spacing}
\startdatakeys
    \source{inconnue}
    \BO{2026}
    \cycle{4}
    \competencesDuSocle{}
    \objectifsApprentissage{B3cb, D3bb, D3bd, A3ee, A3ed}
    \automatismes{C3ab, C2aa, B1ea}
\enddatakeys
\begin{document}
\noindent Dans cet exercice, les questions sont indépendantes.\par\noindent
\textbf{Aucunes justification n'est demandée.}
\begin{enumerate}
    \item \begin{minipage}[t]{.5\linewidth}
        Dans le triangle \mathspoint{ABC}, rectangle en
        \mathspoint{A}, quel calcul doit-on effectuer pour
        déterminer le cosinus de l'angle $\widehat{\mathspoint{ABC}}$~?
    \end{minipage}
    \begin{tikzpicture}[scale=0.75, baseline=(current bounding box.center)]
        \coordinate[label=below right:A](A)at(4,0);
        \coordinate[label=above right:B](B)at(4,3);
        \coordinate[label=below left:C](C)at(0,0);
        \tkzMarkRightAngle[size=0.5,colored](B,A,C)
        \draw(A)--(B)--(C)--cycle;
    \end{tikzpicture}
    \answerspace{1em}
    \item Les notes obtenues par un élève sont 8 ; 12 ; 6 ; 19 ; 15\par
    \begin{enumerate}
        \item Que vaut la médiane de cette série de notes~?
        \answerspace{1em}
        \item Que vaut la moyenne de cette série de notes~?
        \answerspace{1em}
    \end{enumerate}
    \item \begin{minipage}[t]{.5\linewidth}
        Répondre au questions suivantes par lecture graphique (voir ci-contre).
        \begin{enumerate}
            \item Quel est l'antécédent de 2 par la fonction $f$~?
            \answerspace{1em}
            \item Quel est l'image de -1 par la fonction $f$~?
            \answerspace{1em}
            \item Quel est l'image de 0 par la fonction $f$~?
            \answerspace{1em}
            \item Quel est l'antécédent de 0 par la fonction $f$~?
            \answerspace{1em}
        \end{enumerate}
    \end{minipage}\hskip30pt
    \begin{tikzpicture}[scale=0.9, baseline=(current bounding box.north)]
        \simpleFunc xmin=-3,xmax=5,ymin=-4,ymax=4,gridstep=0.5,func=f=-1.5*x+0.5;
    \end{tikzpicture}
    \item Écrire l'expression factorisée de $4x^2-1$.
    \answerspace{1em}
    \item Quelles sont les solutions de l'équation ci-dessous~?
    $$(2x-6)(5+x)$$
    \answerspace{1em}
    \item \begin{minipage}[t]{.5\linewidth}
        La figure ci-contre n'est pas en vrai grandeur.\par\noindent
        Le triangle \mathspoint{DEF} est-il rectangle~?
    \end{minipage}\hskip30pt
    \begin{tikzpicture}[scale=1, baseline=(current bounding box.center)]
        \coordinate[label=below left:A](A)at(0,0);
        \coordinate[label=below right:B](B)at(4,0);
        \coordinate[label=above:C](C)at(4,3);
        \draw(A)--(B)--(C)--cycle;
    \end{tikzpicture}
\end{enumerate}
\end{document}